\documentclass[a4paper]{article}
\usepackage{geometry}
\usepackage{amssymb}
\usepackage{pdfpages}
\usepackage{graphbox}
\usepackage{tikz}
\usepackage{graphicx}  % for \rotatebox{}, used in fisher_table_values.R
\usepackage{multirow}  % for \multirow{},  used in fisher_table_values.R
\begin{document}

\begin{center}
  {\Large\bf Ninety-fifth percentiles of the Fisher distribution}
\end{center}

\rule{0mm}{5mm}
\pagestyle{empty}
\thispagestyle{empty}
\centering
\begin{tabular}{rr|ccccccccc}
  &&\multicolumn{9}{c}{Numerator degrees of freedom\rule[-4mm]{0mm}{4mm}}\\
  & & {\large 1} & {\large 2} & {\large 3} & {\large 4} & {\large 5} & {\large 10} &{\large 15} &{\large 20}& {\large 30}\\[1mm]  \hline
\input{fisher_table_values.txt}   % created by fisher_number_maker.R
  \hline
\end{tabular}

\begin{centering}
\includegraphics[align=c,width=2.2in]{fisher_diagram.pdf}
\end{centering}

\noindent For example, to calculate the $95^\mathrm{th}$ percentile of
the Fisher distribution with 1 and 2 degrees of freedom, use the first
column and the second row to find 18.5.  Thus, if~$X\sim F^1_2$,
then~$P(X\leq 18.5) = 0.95$.

\begin{tikzpicture}[remember picture,overlay]
  \node[anchor=south east,
    xshift=-1cm,
    yshift=1.1cm]
  at (current page.south east){\includegraphics[width=1in]{AUT-logo-block.pdf}};
  \node[anchor=south west,
    xshift=1cm,
    yshift=1.1cm]
  at (current page.south west){{\tt https://github.com/RobinHankin/tables.git}};
\end{tikzpicture}




\end{document}

